\chapter{Methods for preparation and processing}\label{title:processing}
Raw sensor data often contains  noise, inconsistencies, or measurement errors that can negatively affect further analysis and decision-making. Therefore, preprocessing steps are required to clean, transform, and validate the recorded measurements before they are used in algorithms or applications.

This chapter introduces common methods for preparing and processing sensor data. It presents different filtering techniques for noise reduction, explains calibration and normalization procedures, and discusses strategies for handling missing values, outliers, and noisy measurements.

\section{Noise removal and smoothing}\label{title:noise-removal}
Sensor signals often contain high-frequency fluctuations and disturbances from electronics or the environment. Noise removal and smoothing aim to reduce these variations while preserving meaningful signal features. Common approaches include simple statistical filters, such as moving averages, and model-based estimators, like Kalman or orientation filters, which infer the most likely state from the noisy measurement.

The following sections will discuss specific methods, including Kalman filters, Madgwick orientation filters, moving-average and moving-median filters, and radial basis function networks, detailing their operation and typical use cases.

\subsection{Kalman}\label{title:kalman}
The Kalman Filter (a.k.a. Kalman-Bucy-Filter or Stratonovich-Kalman-Bucy-Filter) is a mathematical procedure to determine a dynamic size through imprecise or noisy linear observations. The Kalman Filter has a large usage spectrum and is not just restricted to physics, as the following examples portray this:

\begin{itemize}
    \item Prediction of the position of a moving system, based on acceleration and velocity
    \item Prediction of the rotation of a system, based on torque and angular velocity
    \item Prediction of the remaining duration of a battery, based on previous activity
    \item Numeric weather prediction
    \item Dynamic stochastic general equilibrium (Macroeconomic method for analyzing economic growth and economic situation)
\end{itemize}

\noindent\cite{LRWeb:03, LRWeb:04, LRWeb:05}

{\large \textbf{Fundamental Operating principle}}\\[0.5em]
The Kalman Estimator operates iteratively by decreasing the estimated uncertainty (variance) over time. In the example of predicting the position of an object, the flow might look something like this~\ref{fig:kalman}

\begin{figure}[H]
    \centering
    \includegraphics[width=0.85\linewidth]{images/Kalman-Filter.jpg}
    \caption{Flow of the Kalman Filter}\label{fig:kalman}
\end{figure}
\noindent\cite{LRWeb:07}

{\large \textbf{1. Initialize the state with observation data}}\\[0.5em]
Before being able to predict a value, an initial state is necessary beforehand. Let's take a car for example, we want to predict its current position. The initial state might look like this:
\[
x = \begin{bmatrix}
position\\velocity
\end{bmatrix}
=
\begin{bmatrix}
10\\0
\end{bmatrix}
\]

At the start ($t = 0$), we have an estimated position of ten and a velocity of zero. We can add a covariance matrix which can look like this:

\[
\ P_0 = \begin{bmatrix}
    5 & 0\\
    0 & 1
\end{bmatrix}
\]

A position variance of 5 indicates relatively high uncertainty in our estimate, while a variance of 1 for velocity suggests we are more confident in that measurement.

\noindent {\large \textbf{2. Reinitialize the state}}\\[0.5em]
It is possible that we might realize, that our data is incomplete or unreliable, if this is the case, we can reset our initial data and update it. (This step is optional, thus some applications of the Kalman Filter skip this step.)\\\\
{\large \textbf{3. Predict the state}}\\[0.5em]
The next step is for the filter to predict where the car is next, for the prediction, we have the base formula:

\[
\ \hat{x}_{k|k-1} = F_k\hat{x}_{k-1} + B_k u_k
\]

\[
F_k =
\begin{bmatrix}
1 & \Delta t \\
0 & 1
\end{bmatrix}, \quad
\text{Change of the car's state compared to the previous timestamp.}
\]\\
$\hat{x}_{k|k-1} = $ State of the car at timestamp  $k$ given the information up to timestamp $k - 1$\\
$B_k = $ Control matrix (weights for e.g.\ throttle or acceleration)\\
$u_k = $ Control input (input- e.g.\ throttle or acceleration)\\

In this example we suppose, that the car moves at a constant speed of $1~\frac{m}{s}$ so $B$ is empty, giving us the following formula:\\
\[
\ \hat{x}_{k|k-1} = F_k\hat{x}_{k-1}
\]

The next step is to calculate the new uncertainty:\\
\[
\ P_{k|k-1} = FP_{k-1}F^T + Q
\]

$Q$ being the process noise (e.g.\ wind, friction, etc.)

{\large \textbf{4. Calculate the Kalman Gain}}\\[0.5em]
The Kalman Gain tells us how trustworthy the new observation is, compared to our prediction. To calculate this, we need to substitute into the following formula:

\[
K_k=P_{k|k-1}H^T{(HP_{k|k-1}H^T + R)}^{-1}
\]

$H = $ measurement matrix (Relation from the measurement to the state)\\
$R = $ measurement noise (Noisiness of the sensor)

Hence $R$ and $K$ are reciprocal proportional, with this, two conclusions can be drawn:\\
\small{If the measurement is very accurate $\rightarrow$ $R$ is small $\rightarrow$ $K$ is large $\rightarrow$ we trust the measurement more} \\
\small{If the measurement is very noise $\rightarrow$ $K$ is small $\rightarrow$ $R$ is large $\rightarrow$ we trust our prediction more} \\

Thus $K$ decides, how much our prediction is corrected when new data arrives.

{\large \textbf{5. Calculate the System State and System State Error Covariance Matrix}}\\[0.5em]
The first step here is to calculate the difference between the observation and the predicted value:\\
\[
y_k=z_k-H\hat{x}_{k|k-1}
\]

$z_k = \text{position at timestamp k}$\\

The second step is to update the current state, considering the Kalman Gain:

$\hat{x}_k=\hat{x}_{k|k-1} + K_k y_k$\\

We are multiplying with $K$ (Kalman Gain) here as a way to include the uncertainty.\\

Finally we update the Uncertainty:
\[
P_k = (I - K_k H)P_{k|k-1}
\]

$I = $Identity-Matrix

This leaves us with:
\begin{itemize}
    \item a new improved state
    \item a smaller uncertainty
\end{itemize}

The next operation is to go back to step three and continue from there on again.

\noindent\cite{LRWeb:07, LRWeb:09, LRWeb:10} \\

{\large \textbf{Variations of Kalman}}\\[0.5em]
The Kalman Filter has several forms, the basic functionality is the same but they have some extension, which suit individual scenarios better.

\begin{itemize}
    \item \textbf{Alpha Beta Filter:} Incorporates two parameters: $\alpha$ and $\beta$\\
        $\alpha$: Trustworthiness of the observation\\
        $\beta$: Flexibility of velocity of adaption
    \item \textbf{Extended Kalman Filter:} Observation models need to be non-linear.
    \item \textbf{Unscented Kalman Filter:} Amends the Extended Kalman Filter with Sample Points, solving the approximation issues of its predecessor.
    \item \textbf{Discriminative Kalman Filter:} Uses Bayes' rule to estimate a state based on observations. Especially useful for models whose functions are highly non-linear and/or non-Gaussian.
    \item \textbf{Adaptive Kalman Filter:} Allows the continuous tune of parameters.
\end{itemize}
\noindent\cite{LRWeb:06, LRWeb:03, LRWeb:08}\\

\subsection{Madgwick}\label{title:madgwick}
The Madgwick Orientation Filter is applicable to IMMUs that collect orientation data on all three axes.

Madgwick uses quaternions to represent the orientation of an object, rather than a Euler-representation. With this approach, we can no longer encounter a Gimbal Lock~\ref{title:gimbal} and we have access to a gradient-descent algorithm to calculate the error of the gyroscope measurements as a quaternion derivative.

The technical innovations include:

\begin{itemize}
    \item One customizable parameter defined by system characteristics.
    \item Gyroscope drift compensation
    \item An optimized gradient-descent algorithm with a high performance with a low sampling rate.
\end{itemize}

\noindent\cite{LRWeb:49}

The output of the filter is an estimate of the current orientation of an object. This is notated with a pre-subscript $\overset{W}{I}q$, with $q$ being the orientation in the form of a Quaternion, from the inertial frame $I$ to the world frame $W$.

The idea of Madgwick is to predict $\overset{I}{W}q_{t+1}$ by combining attitude estimates by integrating the measurements taken by the gyroscope $\overset{I}{W}q_\omega$ and the angular velocity acquired by the accelerometer.

The estimation is done by using the following gradient descent algorithm to solve the minimization problem.

\[
\begin{aligned}
&\min_{ {}^{I}_{W}\mathbf{\hat{q}} \in \mathbb{R}^{4 \times 1} } f\left( {}^{I}_{W}\mathbf{\hat{q}}, {}^{W}\mathbf{\hat{g}}, {}^{I}\mathbf{\hat{a}} \right) \\ \\
&f\left( {}^{I}_{W}\mathbf{\hat{q}}, {}^{W}\mathbf{\hat{g}}, {}^{I}\mathbf{\hat{a}} \right) = {}^{I}_{W}\mathbf{\hat{q}}^{*} \otimes {}^{W}\mathbf{\hat{g}} \otimes {}^{I}_{W}\mathbf{\hat{q}} - {}^{I}\mathbf{\hat{a}}
\end{aligned}
\]

$q^*$ denotes the conjugate of $q$
$^W\mathbf{\hat{g}}$ being the normalized gravity vector
$^I\mathbf{\hat{a}}$ being the normalized acceleration measurements

The Madgwick Filter follows three iterative steps.

{\large \textbf{1. Acquire sensor measurements}}\\[0.5em]
Gyroscope data ($^I\mathbf{\omega}_t$) and acceleration measurements ($^I\mathbf{a}_t$) are obtained from the sensor, with $^I\mathbf{\hat{a}}_t$ being the normalized acceleration measurements.

{\large \textbf{2(a). Orientation increment from acceleration}}\\[0.5em]
The next step is to calculate the increment in orientation from the data recorded by the accelerometer.

Madgwick provides an error function:

\[
\ f(^I_W\hat{q}_{est,t}, ^W\hat{g}, ^I\hat{a}_{t+1}) = \begin{bmatrix}
    2 (q_2q_4 - q_1q_3) -a_x \\
    2(q_1q_2 + q_3q_4) - a_y \\
    2 (\frac{1}{2} - q^2_2 - q^2_3) - a_z
\end{bmatrix}
\]

This error function computes the error by subtracting the measured gravity direction from the predicted gravity direction (quaternion)

Next, Madgwick also defines a Jacobian Matrix that tells how small changes in the quaternion affect gravity alignment.

\[
\ J(^I_W\hat{q}_{est,t}, ^W\hat{g}) = \begin{bmatrix}
    -2q_3 & 2q_4 & -2q_1 & 2q_2 \\
    2q_2  & 2q_1 & 2q_4 & 2q_3 \\
    0 & -4q_2 & -4q_3 & 0
\end{bmatrix}
\]

To now compute the orientation, we simply need to multiply the Jacobian Matrix with the error function.

\[
\nabla f = J^T f
\]

The final instruction is to update the term, based on the defined parameter $\beta$

\[
^I_W q_{\nabla, t+1} = -\beta\frac{\nabla f}{||\nabla f||}
\]

{\large \textbf{2(b). Orientation increment from gyroscope}}\\[0.5em]
Now we have to do the exact same for the gyroscope measurements. This is done via a numerical integration

\[
^I_W\dot{q}_{\omega,t+1} = \frac{1}{2} \\ ^I_W\hat{q}_{est,t} \otimes{{[0, ^I\omega_{t+1}]}^T}
\]

{\large \textbf{3. Fuse measurements}}\\[0.5em]
The final step is to fuse the acquired data from both the accelerometer and the gyroscope, to obtain an estimate of the attitude

\[
^I_W q_{est,t+1} = ^I_W\hat{q}_{est,t} + ^I_W\dot{q}_{est,t+1} \\ \Delta t
\]

$\Delta t$ being the time between the two samples $t$ and $t + 1$

The next operation is to go back to step one and continue the process from there on again.

The customizable parameter $\beta$ defines how much the accelerometer corrects the gyroscope. If $\beta$ is large, the accelerometer has more affect than the gyroscope, causing a fast correction but also causing noise. On the other hand, a small $\beta$ is smoother but causes more drift. 

Initially, the estimate of the orientation, biases and the sampling times have to be passed to the filter, to work correctly. The bias is calculated by taking the average/mean of the samples with a resting IMU\@.

\noindent\cite{LRWeb:50}

\subsection{Moving-Average-Filter}\label{title:maf}
The arithmetic average is determined by the function:

\[
\bar{x} = \frac{1}{n} \sum_{i=1}^{n} x_i
\]

When trying to smooth out noise with time-series data, the arithmetic average values data further in the past, the exact same as data which has just been recorded. The lack of distinction between the recording of values has several problems:

\begin{itemize}
    \item Indefinitely previous data barely affects future values
    \item Heavy outliers have a large effect on the result
    \item For large datasets, new values barely affect the average
\end{itemize}

In order to solve this issue, we need a method which allows us to time-box the values used to calculate the average. To achieve this, we use the Moving-Average Filter (MAF)

\begin{quote}
    The moving average filter computes the average value of a set of samples within a predetermined window.
\end{quote}
\noindent\cite{LRWeb:01}

The MAF holds the index of a current position, a window size, and an array of values. The opening value is determined by taking the initial average of the first n values. The next averages are calculated by moving the position forward. The graph of a noisy function, compared to the same function smoothed with the MAF is portrayed in figure~\ref{fig:mfa}.

\begin{figure}[H]
    \centering
    \includegraphics[width=0.55\linewidth]{images/MFA-Graph.png}
    \caption{Graph of a function smoothed with the Moving-Average-Filter}\label{fig:mfa}
\end{figure}
\noindent\cite{LRWeb:02}

\subsection{Moving Median Filter}\label{title:mmf}
The median of an ordered data set $x_1, x_2, \ldots, x_n$ is defined as

\[
\text{Median}(x) =
\begin{cases}
x_{\frac{n+1}{2}}, & \text{if } n \text{ is odd} \\[6pt]
\dfrac{x_{n/2} + x_{n/2 + 1}}{2}, & \text{if } n \text{ is even}
\end{cases}
\]


The Moving Median Filter works the exact same way as the Moving Average Filter, the only twist being, that we use the median instead of the mean for calculating the average in a given window. This procedure has the following characteristics compared to the MAF~\ref{tab:moving-comparison}:\\

\begin{table}[ht]
\centering
\begin{tabular}{lcc}
\toprule
\textbf{Characteristic} & \textbf{Moving Average Filter} & \textbf{Moving Median Filter}\\
\midrule
Speed                       & High      & High --- Medium \\
Complexity                  & $O(n)$      & $O(n \cdot \log n)$    \\
Outlier affection           & High      & Low           \\
Implementation complexity   & Simple    & Moderate      \\
\bottomrule
\end{tabular}
\caption{Comparing the Moving Average to the Moving Median}\label{tab:moving-comparison}
\end{table}

The higher computation cost of the MMF is due to the fact, that the window of values have to be sorted every time when estimating.\\
\noindent\cite{LRWeb:11}

\subsection{Radial Basis Function Network}
\begin{quote}
    A radial basis function network is an artificial network whose activation function is a radial basis function.
\end{quote}
\noindent\cite{LRWeb:15} 

\begin{figure}[H]
    \centering
    \includegraphics[width=0.75\linewidth]{images/Radial-Basis-Function-Network.png}
    \caption{Radial Basis Function Network (source: own work)}\label{fig:rbfn-example}
\end{figure}
\noindent\cite{LRWeb:16}

The Radial Basis Function Network (RBFN) utilizes three layers in its functionality, one input-layer, one hidden layer, and an output-layer.

The input layer consists of only one vector that should be classified. The entire vector is passed to the RBF neurons.\\
Every neuron consists of a prototype vector\footnote{The prototype vector is also called the center of the neuron}, which the neuron ``learned'' from the training. When a neuron receives the input vector, it compares the input with its prototype, displaying a number in the range $[0, 1]$, this value is called the \textbf{activation value}.\\

\[
\ \phi(x) = e^{-\beta||x-\mu||^2}
\]

$x$ = Input vector\\
$\mu$ = Prototype vector\\
$\beta$ = Height of Gaussian (learned by nodes during training)\\

The values mean:

\[
\ ||x-\mu|| \to \infty \Rightarrow \text{Input vector is very different from the prototype vector}
\]
\[
\ \phi(x) = 1 \Rightarrow \text{Input vector is identical to the prototype vector}
\]

The output layer consists of nodes (one for every class we are trying to categorize), which calculates a result for its category. The result is determined by taking the weight value for a RBF-Neuron\footnote{The weight is estimated by the output nodes} and multiplying it with the activation value. The final output is a score for every class, where usually the highest value is used to classify the input vector.

\noindent\cite{LRWeb:16}

{\large \textbf{Usage}}\\[0.5em]
The RBFN is usually used for classification problems whose data points cannot be separated by a linear function. A graph, for which this is true, can look like figure~\ref{fig:rbfn}.

\begin{figure}[H]
    \centering
    \includegraphics[width=0.5\linewidth]{images/RBFN-Example.png}
    \caption{Graph of data points, which can't be separated by a linear function}\label{fig:rbfn}
\end{figure}

A field of usage can be the correction of Sensor Drift (see Section~\ref{title:drift}). There should also be an alarm policy, which handles the Drift. If Drift is detected, then act. Actions can vary from solutions based on Artificial Intelligence (Neural Network or other Regression-models trained on previous data) to algorithmic solutions (Kalman Filter variation, see Section~\ref{fig:kalman}).

\noindent\cite{LRWeb:16}

\section{Calibration, Scaling and Normalization}
\boldtitle{Calibration}
\begin{quote}
    In scientific terms, calibration is defined as the process of configuring an instrument to provide a result for a sample within an acceptable range. 
\end{quote}
\noindent\cite{LRWeb:25}

Sensor calibration is a way to determine whether the values a sensor records are faulty. The main advantages of a calibrated sensor include:
\begin{itemize}
    \item Consistency
    \item Improved performance
    \item More accurate data
\end{itemize}

The calibration process depends on several factors such as the environment, sensor type, and required accuracy. Common calibration methods include:
\begin{itemize}
    \item \textbf{Manual calibration:} Comparing the sensor output to a known reference and adjusting it accordingly.
    \item \textbf{Automated calibration:} Software-driven calibration, where the software adjusts the sensor settings based on a reference standard.
    \item \textbf{Two-Point calibration:} Calibration is performed at two reference points, the zero point (no input) and the span point (a known full-scale input.\footnote{Two-Point Calibration is commonly used for sensors that record linear values such as temperature or pressure sensors}).
    \item \textbf{Multi-Point calibration:} The sensor is calibrated at multiple points to ensure accuracy across the entire measurement range.
\end{itemize}

Regardless of the chosen method, sensors require regular calibration to ensure data reliability over time.\ \cite{LRWeb:25}

\boldtitle{Scaling}
Scaling is the process of converting a recorded sensor value's unit (which might be different from the unit needed) to a desired unit. This is a necessary step and cannot be skipped, as it is essential for the algorithms to work with the right unit. Examples for Scaling are:
\begin{itemize}
    \item The Madgwick Filter~\ref{title:madgwick} ($^\circ/s \rightarrow rad/s$)
    \item Temperature Sensors ($^\circ\mathrm{K} \rightarrow ^\circ\mathrm{C}$)
    \item Accelerometers ($km/h² \rightarrow m/s²$)
\end{itemize}
\noindent\cite{LRWeb:22}

\boldtitle{Normalization}
Normalization adjusts values to a common scale, making data from multiple sensors comparable and easier to process. Many sensor fusion algorithms, such as the Madgwick~\ref{title:madgwick} or Mahony filters, require normalized input data.

One common normalization technique is \textbf{min-max feature scaling}, which maps data to the interval $[0; 1]$ using the formula:

\[
x' = \frac{x - x_{\min}}{x_{\max} - x_{\min}}
\]

Another widely used approach is the \textbf{Z-score normalization method}~\ref{title:zscore}, which scales values based on the mean and standard deviation of the dataset.\\
\noindent\cite{LRWeb:23, LRWeb:24}

\section{Handling missing measurements, noisy values and outliers}
When working with sensor data, inconsistent measurements should not be part of the final usage. They can either be normalized to fit the other data (if possible) or they can be dropped all together.

\boldtitle{Missing measurements}
Missing measurements can happen for a variety of reasons; some of them are as follows:

\begin{itemize}
    \item Sensor outage
    \item Communication issues
    \item Data corruption or data deletion
\end{itemize}

Single missing values might not cause a lot of trouble, but if data is missing for a long time, systems working with those measurements will not work until new observations are received.

There are multiple possibilities to fight this issue. Simple methods might be to just use zero values for the missing data, a more refined way is to use some sort of averaging in consideration of the time; good contenders therefore might be the Moving Average Filter~\ref{title:maf} or the Moving Median Filter~\ref{title:mmf}.

A more complex sophisticated approach is to use methods like interpolation, regression or even machine learning system, although they might be overkill in some scenarios.

\noindent\cite{LRWeb:43}

\boldtitle{Noise in data}
Noise in measurements can usually be fixed by occasionally calibrating the sensor to fit its environment again. It is also important to test the sensors and to maintain said sensors, as e.g.\ dirt or mud (applicable for EMUs) can negatively affect sensors by adding unnecessary noise.

\noindent\cite{LRWeb:44}

\boldtitle{Outliers}
After we detected an outlier in our observations~\ref{title:outlier}, there are multiple approaches to deal with them. The simplest form is to just remove the outliers completely. In some use cases this might not be optimal as we are losing information. Another possibility is to clamp it (aka winsorize), meaning we have predetermined values for the minimum and maximum and if a measurement exceeds either, they will just be treated as the minimum/maximum. Ultimately it is also possible to treat outlier values as null values, so that they don't affect the calculations.

\noindent\cite{LRWeb:42}